% Malla curricular · Filosofía 10.° · Periodo IV · PILOTO (ejemplo completo del kit)
% Temas: los vistos en el periodo IV de 2026 (ver evaluaciones/ejemplos/2026-11-14_filosofia-10.tex).
% Referente: MEN, Documento 14 (2010). Enfoque del PEI: Escuela Activa.
% Compilar: bash curriculo/plantilla/compilar.sh curriculo/filosofia/10/malla-P4.tex
\documentclass[malla]{planeacion}
\materia{filosofia}
\grado{10.°} \periodo{IV} \anio{2027}
\begin{document}
\begin{malla}{¿Pueden convivir la fe y la razón? ¿Qué cambió en la manera de pensar al ser humano
  cuando Europa pasó de la escolástica medieval al humanismo del Renacimiento?}
  \fila{\textbf{1. La herencia griega en la Edad Media.} Platón y Aristóteles; cómo llegaron sus
    obras a la Europa medieval: neoplatónicos, traducciones del árabe y del griego (escuela de
    Toledo), Averroes. {\scriptsize[FIL-10-P4-D1]}}
    {Conocimiento}
    {¿Cómo se transmiten las ideas a lo largo de la historia y quién las conserva?}
    {\textbf{Crítica:} sitúa las ideas en su contexto histórico y reconoce de dónde provienen.}
    {\punto Ubica en una línea de tiempo a Platón, Aristóteles, Agustín y Tomás y calcula las
     distancias entre ellos.
     \punto Interpreta \emph{La escuela de Atenas} (Rafael) como síntesis de dos maneras de
     filosofar.}
  \fila{\textbf{2. San Agustín de Hipona (354--430).} La interioridad («vuelve a ti mismo»), la
    inquietud del corazón (\emph{Confesiones} I, 1), fe y razón: «comprende para creer, cree
    para comprender». {\scriptsize[FIL-10-P4-D2]}}
    {Conocimiento y moral}
    {¿Dónde busca el ser humano la verdad: fuera o dentro de sí?}
    {\textbf{Dialógica:} interpreta textos filosóficos breves y los pone en diálogo con su
     experiencia.}
    {\punto Identifica la tesis de un fragmento de Agustín.
     \punto Explica la relación entre fe y razón que propone y la compara con una postura
     distinta.}
  \fila{\textbf{3. Tomás de Aquino (1225--1274) y la escolástica.} Las universidades medievales;
    la recuperación de Aristóteles; el método de la \emph{quaestio} (objeciones, «en sentido
    contrario», respuesta, solución de las objeciones) en la \emph{Suma teológica}.
    {\scriptsize[FIL-10-P4-D3]}}
    {Conocimiento}
    {¿Cómo se argumenta con rigor frente a quien piensa distinto?}
    {\textbf{Crítica y creativa:} reconstruye argumentos y aplica un método de argumentación a
     un problema nuevo.}
    {\punto Explica la lógica de la \emph{quaestio}.
     \punto Formula una \emph{quaestio} escolástica sobre un problema actual, con objeciones y
     respuesta.}
  \fila{\textbf{4. ¿Fue «oscura» la Edad Media?} El origen de la etiqueta; Renacimiento, imprenta
    y difusión de las ideas; el humanismo y Erasmo de Rotterdam (\emph{Elogio de la locura},
    1511). {\scriptsize[FIL-10-P4-D4]}}
    {Conocimiento (historia de las ideas)}
    {¿Qué cambia cuando cambia la forma de transmitir el conocimiento?}
    {\textbf{Crítica:} evalúa una afirmación histórica con evidencias a favor y en contra.}
    {\punto Evalúa la etiqueta «Edad Oscura» con evidencias.
     \punto Relaciona la imprenta con la difusión de las ideas a partir de datos.
     \punto Caracteriza el humanismo frente a la escolástica.}
  \fila{\textbf{5. Tomás Moro y la \emph{Utopía} (1516).} Crítica a la sociedad inglesa (las
    «ovejas que devoran hombres»), la sociedad ideal de la isla de Utopía, sentido del término.
    {\scriptsize[FIL-10-P4-D5]}}
    {Moral y política}
    {¿Es útil imaginar una sociedad perfecta?}
    {\textbf{Creativa y crítica:} imagina alternativas a la realidad social y evalúa su
     viabilidad y sus riesgos.}
    {\punto Interpreta la metáfora de las ovejas como crítica social.
     \punto Infiere la intención de una medida de la sociedad utópica.
     \punto Propone y defiende una «utopía» propia para su comunidad.}
\end{malla}
\end{document}
